\documentclass[a4paper,11pt]{article}

\newcommand{\TPnumero}{Lab 4 -- Unit roots and ARIMA}
% =========================================================
% Tutorial / lab preamble — Time Series Analysis module
% article class + fancyhdr + tcolorbox + listings (English)
% Author : Aymen Ben Brik (aymen.benbrik@esprit.tn)
% =========================================================

\usepackage[utf8]{inputenc}
\usepackage[T1]{fontenc}
\usepackage[english]{babel}
\usepackage{lmodern}
\usepackage{amsmath, amssymb, amsthm, mathtools, bm}
\usepackage{graphicx}
\usepackage{booktabs}
\usepackage{array}
\usepackage{multirow}
\usepackage{colortbl}
\usepackage{xcolor}
\usepackage{tikz}
\usetikzlibrary{positioning, arrows.meta, calc, shapes, fit, backgrounds, matrix}
\usepackage{listings}
\usepackage[most]{tcolorbox}
\usepackage{geometry}
\geometry{a4paper, margin=2cm, headheight=15pt}
\usepackage{fancyhdr}
\usepackage[hidelinks]{hyperref}
\usepackage{enumitem}

% --- Colours ---
\definecolor{espritBlue}{RGB}{30, 60, 110}
\definecolor{espritAccent}{RGB}{200, 80, 30}
\definecolor{boxEx}{RGB}{30, 60, 110}
\definecolor{boxQ}{RGB}{180, 120, 0}
\definecolor{boxC}{RGB}{30, 100, 60}
\definecolor{boxAtt}{RGB}{200, 80, 30}

% --- Listings ---
\definecolor{codebg}{RGB}{248,248,248}
\definecolor{codeframe}{RGB}{220,220,220}
\definecolor{keyword}{RGB}{0,82,155}
\definecolor{comment}{RGB}{88,110,117}
\definecolor{string}{RGB}{133,153,0}

\lstdefinestyle{pythonstyle}{
  language=Python,
  basicstyle=\ttfamily\small,
  keywordstyle=\color{keyword}\bfseries,
  commentstyle=\color{comment}\itshape,
  stringstyle=\color{string},
  backgroundcolor=\color{codebg},
  frame=single,
  rulecolor=\color{codeframe},
  frameround=tttt,
  breaklines=true,
  showstringspaces=false,
  tabsize=2
}
\lstdefinestyle{Rstyle}{
  language=R,
  basicstyle=\ttfamily\small,
  keywordstyle=\color{keyword}\bfseries,
  commentstyle=\color{comment}\itshape,
  stringstyle=\color{string},
  backgroundcolor=\color{codebg},
  frame=single,
  rulecolor=\color{codeframe},
  frameround=tttt,
  breaklines=true,
  showstringspaces=false,
  tabsize=2
}
\lstset{style=pythonstyle}

% --- Common macros (configurable path) ---
\providecommand{\macrosPath}{../../preamble/macros.tex}
\input{\macrosPath}

% --- Exercise counter ---
\newcounter{excounter}
\renewcommand{\theexcounter}{\arabic{excounter}}

\newtcolorbox[use counter=excounter]{exercise}[1][]{
  enhanced, breakable,
  colback=boxEx!5, colframe=boxEx, fonttitle=\bfseries,
  title={Exercise~\theexcounter\ifstrempty{#1}{}{ -- #1}},
  arc=2pt, boxrule=0.6pt
}
\newtcolorbox{question}[1][]{
  enhanced, breakable,
  colback=boxQ!5, colframe=boxQ, fonttitle=\bfseries,
  title={Question\ifstrempty{#1}{}{ -- #1}},
  arc=2pt, boxrule=0.5pt
}
\newtcolorbox{solution}[1][]{
  enhanced, breakable,
  colback=boxC!5, colframe=boxC, fonttitle=\bfseries,
  title={Solution\ifstrempty{#1}{}{ -- #1}},
  arc=2pt, boxrule=0.5pt
}
\newtcolorbox{warning}[1][]{
  enhanced, breakable,
  colback=boxAtt!5, colframe=boxAtt, fonttitle=\bfseries,
  title={Warning\ifstrempty{#1}{}{ -- #1}},
  arc=2pt, boxrule=0.5pt
}

% --- Headers / footers ---
\pagestyle{fancy}
\fancyhf{}
\fancyhead[L]{\textsc{\moduleNom} -- \auteurNom}
\fancyhead[R]{\TPnumero}
\fancyfoot[C]{\thepage}
\fancyfoot[L]{\auteurInstitution}
\fancyfoot[R]{\auteurEmail}
\renewcommand{\headrulewidth}{0.4pt}
\renewcommand{\footrulewidth}{0.2pt}

% --- Placeholder ---
\providecommand{\TPnumero}{Lab}


\title{Lab 4: Unit-root tests and ARIMA modelling\\
\large ADF, KPSS, differencing, and ARIMA fitting on real data}
\author{\auteurNom\\\auteurEmail\\\auteurInstitution}
\date{\anneeUniversitaire}

\begin{document}
\maketitle

\section*{Goals}
\begin{itemize}
  \item Distinguish a trend-stationary (TS) process from a
        difference-stationary (DS) process via simulation.
  \item Run the Augmented Dickey--Fuller (ADF) and KPSS tests, read
        their $p$-values, and combine them.
  \item Apply differencing to remove a unit root and verify the
        transformation by re-running the tests.
  \item Fit an ARIMA$(p,d,q)$ model on a non-stationary series and
        produce a short forecast.
\end{itemize}

\section*{Setup}
\begin{itemize}
  \item Python 3 with \texttt{numpy}, \texttt{matplotlib},
        \texttt{statsmodels}.
  \item \texttt{numpy.random.seed(42)}.
  \item Dataset: \texttt{data/SouvenirSales.csv} (Australian souvenir
        shop, monthly, Jan 1995 -- Dec 2001, $T = 84$).
\end{itemize}

\begin{warning}
\texttt{adfuller(y)} returns the tuple
\texttt{(stat, pvalue, usedlag, nobs, critvalues, icbest)}.
\texttt{kpss(y, regression="c")} returns
\texttt{(stat, pvalue, lags, critvalues)} (with a Future\-Warning
about $p$-value clamping that you can ignore).
\end{warning}

% =========================================================
\begin{exercise}[TS vs.\ DS by simulation]
We compare two visually similar series.

\textbf{(1)} Simulate $T = 200$ observations of a trend-stationary
process
\[
  Y_t = 0.05\, t + \varepsilon_t, \qquad
  \varepsilon_t \iid \mathcal N(0, 1).
\]

\textbf{(2)} Simulate $T = 200$ observations of a random walk with drift:
\[
  Z_t = Z_{t-1} + 0.05 + \eta_t, \qquad
  \eta_t \iid \mathcal N(0, 1), \quad Z_0 = 0.
\]

\textbf{(3)} Plot $Y_t$ and $Z_t$ on the same figure. Visually, can
you tell which is TS and which is DS?

\textbf{(4)} Detrend each series by OLS regression on $t$ and plot
the residuals. Comment: the TS residual should look like white
noise; the DS residual still wanders.
\end{exercise}

% =========================================================
\begin{exercise}[ADF test]
\textbf{(1)} Run \texttt{adfuller} on $Y_t$ and on $Z_t$ from
Exercise~1. Report for each series:
\begin{itemize}
  \item the ADF statistic $t_{\widehat\gamma}$;
  \item the $1\%$, $5\%$, $10\%$ critical values;
  \item the $p$-value;
  \item the conclusion at the $5\%$ level (reject $H_0$: unit root,
        or fail to reject).
\end{itemize}

\textbf{(2)} On the random walk with drift $Z_t$, the ADF test should
\emph{fail to reject} the unit-root null. On the trend-stationary
$Y_t$, you must use \texttt{regression="ct"} (constant + trend);
otherwise the deterministic trend would inflate the test statistic.

\textbf{(3)} Repeat the ADF test on $\Delta Z_t = Z_t - Z_{t-1}$.
You should now \textbf{reject} the unit root: differencing has made
the series stationary.
\end{exercise}

% =========================================================
\begin{exercise}[KPSS test]
The KPSS test reverses the null and the alternative:
$H_0$ is stationarity, $H_1$ is unit root.

\textbf{(1)} Run \texttt{kpss(y, regression="c")} on $Y_t$ (after
removing the linear trend, use \texttt{regression="ct"} on the raw
series) and on $Z_t$. Report KPSS statistic, $p$-value, conclusion.

\textbf{(2)} Build the ADF/KPSS combination matrix on the real
series of Exercise~5: which of the four cells (stationary,
trend-stationary, difference-stationary, ambiguous) does it land in?
\end{exercise}

% =========================================================
\begin{exercise}[Differencing the random walk]
\textbf{(1)} Compute $\Delta Z_t = Z_t - Z_{t-1}$ on the simulated
random walk of Exercise 1. Plot $\Delta Z_t$ vs.\ $t$.

\textbf{(2)} Plot the sample ACF of $Z_t$ and of $\Delta Z_t$ up
to lag 20. Compare: the ACF of $Z_t$ decays slowly (signature of
a unit root); the ACF of $\Delta Z_t$ should be near zero for
all $h \geq 1$.

\textbf{(3)} Run \texttt{adfuller} and \texttt{kpss} on $\Delta Z_t$.
Confirm that one differencing was enough.
\end{exercise}

% =========================================================
\begin{exercise}[ARIMA on Souvenir sales]
\textbf{(1)} Load \texttt{data/SouvenirSales.csv} and take logs:
$y_t = \log(\text{sales}_t)$. Plot $y_t$.

\textbf{(2)} Run ADF and KPSS on $y_t$ (with \texttt{regression="ct"}
since there is a clear linear trend on the log scale). Conclude.

\textbf{(3)} Fit the model
\[
  \text{SARIMA}(0, 1, 1)(0, 1, 1)_{12}
\]
on $y_t$ using \texttt{statsmodels.tsa.statespace.sarimax.SARIMAX}.
Report the estimated coefficients and the AIC.

\textbf{(4)} Plot the residuals of the fitted model. Run a Ljung--Box
test at lag 24 on the residuals: are they white noise?

\textbf{(5)} Forecast the next 12 months. Plot the in-sample fit and
the forecast with $95\%$ confidence intervals.
\end{exercise}

\bigskip
\begin{warning}
The KPSS test in \texttt{statsmodels} returns a $p$-value clamped to
$[0.01, 0.10]$. To get a finer reading, use the printed critical
values (1\%, 2.5\%, 5\%, 10\%) directly.
\end{warning}

\end{document}
